\documentclass[12pt]{article}%[11pt]{report}{book}
\usepackage{graphicx,amssymb,latexsym}
%\def\thesection{\Roman{section}}
%\textwidth=6in
%\advance\hoffset -2.5truein
%\renewcommand{\theequation}{\thesection.\arabic{equation}} %set counter at
%\beginning of each section`\setcounter{equation}{0}'
%\textheight %text height in inches
%\topmargin % 0 is the default which gives a one inch margin
%\headheight % about 1/4 in is the default, 0 removes
%\headsep % the distance between the head and the body of the text the
%default is about 1/4 in
%\pagestyle{normal} %{empty}
\def\negate#1{\mbox{$\;$\begin{picture}(10,12)
\put(0,-2){\line(1,1){10}}
\put(6,0){\makebox(0,0)[b]{#1}}
\end{picture}$\;$}}
%\advance\hoffset by .5truein\relax%controls horizontal position on the page
\input bigc
\newtheorem{example}{Example}[section]
\newtheorem{corollary}{Corollary}[section]
\newtheorem{lemma}{Lemma}[section]
\newtheorem{theorem}{Theorem}[section]
\newtheorem{algorithm}{Algorithm}[section]
\newtheorem{proposition}{Proposition}[section]
\newtheorem{definition}{Definition}[section]
\newtheorem{notations}{Notations}[section]
\newtheorem{assumption}{Assumption}[section]
\newtheorem{remark}{Remark}
\newtheorem{problem}{Problem}[section]
\def\span{\mbox{span}}
\begin{document}
\baselineskip=18pt
\begin{center}
{\large \bf A CONSTRUCTIVE METHOD FOR NUMERICALLY COMPUTING CONFORMAL
MAPPINGS FOR GEARLIKE DOMAINS}
\end{center}
\vspace*{.5in}


\noindent {\bf INTRODUCTION}\vspace*{12pt}

The Riemann mapping theorem asserts that the 
open unit disk  ${\cal D}  = \{z | \  \  |z| < 1\}$ is conformally equivalent to 
each simply connected domain  ${\cal G}$  in the complex plane, whose 
boundary consists of at least two points, i.e., there exists a 
function $f$, analytic and univalent function on  ${\cal D}$ , such that $f$ 
maps  ${\cal D}$  onto  ${\cal G}$ .  More precisely, if $d_o$  is an
arbitrary point in  ${\cal D}$  
and $g_o$  is an arbitrary point in  ${\cal G}$, then the Riemann mapping 
theorem asserts that there exists a unique conformal mapping $f$ of 
 ${\cal D}$  onto  ${\cal G}$  such that $f(d_o) = g_o  \  {\rm and} \  f'(d_o) > 0.$  If the boundary of 
 ${\cal G}$  is piece-wise analytic and $g_1$  is a point on the
boundary of  ${\cal G}$, 
then the uniqueness assertion of the Riemann mapping theorem can 
be reformulated alternately as the statement that there exists a
unique conformal  mapping $f$ of  ${\cal D}$ onto  ${\cal G}$ such
that $f(d_o) = g_o \   {\rm and} \  f(1) = g_1 .$

The problem of constructing the explicit conformal mapping 
guaranteed by the Riemann mapping theorem is usually difficult, 
even numerically.  There are constructive proofs \cite{6} of the 
Riemann mapping theorem, but, because of their general nature, 
they often converge only slowly to the desired solution.

For polygonal domains, i.e., simply connected domains bounded by
segments of straight lines, there is a well-known representation
formula, called the Schwarz-Christoffel formula (see \cite{7}, p.
189), which the conformal mapping functions must satisfy.  The
difficulty with using the Schwarz-Christoffel formula to construct
the conformal mapping for an explicit polygonal domain has always
been that the formula contains unknown parameters, called the
accessory parameters, which have to be determined before the mapping
function can be computed.  D. Gaier in his monograph \cite{Ga} on
conformal mapping surveyed several methods which had been proposed for
solving the accessory parameter problem and proposed a method for
reducing the problem to a system of nonlinear
equations.  In 1980 L. N. Trefethen \cite{10} devised an effective
procedure for solving the accessory parameter problem and, hence,
constructing the conformal mapping function for a polygonal domain.
Trefethen used the geometry of the polygonal domain to construct a
constrained system of nonlinear equations, the solution set of which
would be the accessory parameter set.

In \cite{3}, following work of A. W. Goodman \cite{5}, necessary
and sufficient conditions were given for representing the conformal
mapping function which would map ${\cal D}$ onto a image domain
${\cal G}$ which is gearlike.  A gearlike domain is a simply
connected domain which contains the origin and is bounded by arcs of
circles centered at the origin and by segments of straight lines
through the origin.  (See Figure 1.)  We note that the logarithmic
image of a gearlike domain is a periodic polygonal domain.  Hence,
one possible method for solving the conformal mapping problem for a
gearlike domain would be to apply an adaption of Trefethen's method
to the logarithmic image.  In this paper we propose to directly use
the representation formula for gearlike functions and a development
motivated by the work of Trefethen to construct a procedure for
numerically computing the conformal mapping function for a gearlike
domain.  The kernel of that procedure is an algorithm which
prescribes a system of nonlinear equations whose solution set is the
unknown accessory parameter set for the gearlike domain being
considered.  \newpage

$$\includegraphics[scale=0.25]{fig-01.eps}$$


\centerline{Figure 1}
\vspace*{12pt}

We will give several examples using this procedure for constructing
both the forward mappings from ${\cal D}$ to explicit gearlike
domains ${\cal G}_i$
and the inverse mappings from the domains ${\cal G}_i$ back to ${\cal
D}$.  We will
conclude with an application which shows that a coefficient
conjecture of R. W. Barnard [Ba] fails.

There are several recent references which give alternate
constructions for related conformal mapping problems.  K. P. Jackson
and J. C. Mason \cite{6.5} treated a problem of crack stress
around holes in two-dimensional plates by mapping the crack region
conformally to the exterior to the unit circle.  They used an
adaption, made by Trefethen \cite{11}, of the SCPACK programs
Trefethen developed for solving Schwarz-Christoffel mapping problems.
 Recently, P. Bj\o rstad and E. Grosse \cite{3.2} devised a method
for constructing the conformal mapping functions (from the unit disk)
to circular arc polygons.  A circular arc polygon is a polygon where
the sides of the polygon are allowed to be general arcs of circles.
Their method applies an o.d.e. solver to a specific second order
differential equation which represents the circular arc polygon
mapping function.\vspace{12pt}

\noindent {\bf GEARLIKE DOMAINS}\vspace*{12pt}

Let  ${\cal G}$  be a gearlike domain in the complex 
plane with $n$ sides.  Let $w_k , 1 \leq k \leq n,$ denote the vertices of 
 ${\cal G}$, some of which may lie at infinity.  For convenience, we
will also denote $w_n$ as $w_0$.  For each finite vertex $w_k$ let
$\pi \alpha_k$ denote the interior angle at $w_k$ and let
$\pi\beta_k$ denote the exterior angle at $w_k$.  By definition
$\alpha_k$ and $\beta_k$ satisfy the relation $\alpha_k + \beta_k =
1$.  For each infinite vertex $w_k$ set $\beta_k = 1$.  We note, in
this latter case, that $\pi\beta_k$ is generally not the exterior
angle at $w_k$.

For example, in Figure 2 at the corner vertices $w_3$  and
$w_6$  we  
have  $\beta_3  = \frac{1}{2}$ and $\beta_6  =-\frac{1}{2}.$  At
$w_2,$ the tip of a slit, the external angle is $-\pi,$ so  $\beta_2
= -1.$\vspace*{12pt}

$$\includegraphics[scale=0.25]{fig-02.eps}$$

\centerline{Figure 2}
\vspace*{12pt}

In \cite{3} we showed that the following representation formula
describes the conformal function $f$ mapping ${\cal D}$ onto ${\cal
G}$ with $f(0) =
0$\begin{equation}\frac{zf'(z)}{f(z)}=\prod^n_{k=1}(1-\bar{z}_kz)^{-\beta_k}\label{eq:1}\end{equation}
where the points $z_k$ are the preimages on $|z| = 1$ of the vertices
$w_k$ and where the following constraint conditions must be
satisfied: \renewcommand{\theequation}{\roman{equation}} \setcounter{equation}{0}\begin{equation}
\beta_k\in\left\{-1,-\frac{1}{2},0,\frac{1}{2},1\right\}, \  \ 1\leq
k \leq
n,\label{eq:i}\end{equation}\begin{equation}\sum^n_{k=1}\beta_k=0,\label{eq:ii}\end{equation} and \begin{equation}
\sum^n_{k=1}\beta_k Arg(z_k)\equiv 0 \  \ (mod \
\pi).\label{eq:iii}\end{equation}

We note that if ${\cal G}$ is gearlike, then it is easily seen 
for the mapping function $f$ we must have that $zf^\prime (z)/f(z)$ is
either pure real or pure imaginary on the boundary of ${\cal D}$.  If
a function $f$ satisfies (i) and (ii), but not (iii), then $zf^\prime
(z)/f(z)$ will map ${\cal D}$ to a domain bounded by radial segments
which emanate from the origin, but which are twisted away from the
coordinate axes and, hence, $f$ will map ${\cal D}$ to a domain which
is not gearlike.

Equation (\ref{eq:1}) can be solved to explicitly represent the conformal 
mapping $f$ on  ${\cal D}$ as
\renewcommand{\theequation}{\arabic{equation}}\setcounter{equation}{1}
\begin{equation} f(z)=cz
\exp\left\{\int^z_0\frac{\left[\displaystyle{\prod^n_{k=1}}(1-\bar{z}_kw)^{-\beta_k}-1\right]}{w}dw\right\}.\label{eq:2}\end{equation}

In the representation formula (\ref{eq:2}) the parameter $c$ and the 
prevertices $z_k, 1 \leq k \leq n,$ are accessory parameters which must 
be determined before the conformal mapping $f$ can be computed.  
For each accessory parameter set satisfying (i), (ii) and (iii) a  
function $f$ given by (\ref{eq:2}) will map  ${\cal D}$ onto a gearlike
domain $f({\cal D})$  
whose local boundary at each vertex $w(k) = f(z_k)$ will be 
conformally homotopic to the local boundary of ${\cal G}$ at $w_k.$  However, 
corresponding sides of $f({\cal D})$ and  ${\cal G}$  which are arcs of
circles will  
generally not have the same modulus and corresponding sides of
$f({\cal D})$ 
and  ${\cal G}$  which are segments of straight lines will generally not have 
the same argument.  

Using the second formulation for uniqueness in the statement 
of the Riemann mapping theorem we may suppose that the prevertex 
$z_n  = 1,$ i.e.,  that $f(1) = w_n,$ and that $w_n$  is finite.  The later 
assumption can always be made via a suitable reindexing of the 
vertices of  ${\cal G}$.  The parameter $c$ can be computed directly.  If we 
let $$w^*=
\exp\left\{\int^1_0\frac{\left[\displaystyle{\prod^n_{k=1}}(1-\bar{z}_kw)^{-\beta_k}-1\right]}
{w}dw\right\}$$ then $$c = w_n /w^* .$$

Thus, there remain $n-1$ unknown accessory parameters to be
calculated.  To determine the accessory parameters we will construct
a system of $n-1$ real equations, dependent on the parameters
$z_1,\ldots,z_{n-1}$, whose solution set will be the correct
accessory parameter set.

If the domain ${\cal G}$ is starlike, then there exists a unique
function $f$ which will map ${\cal D}$ onto ${\cal G}$ as determined
by the vertices $\{w_k\}$ and the exterior angles $\{\beta_k \pi\}$.
However, if the domain ${\cal G}$ is non-starlike,there exist several
functions which will map ${\cal D}$ onto a domain determined only by
specifying the vertices $\{w_k\}$ and the exterior angles $\{\beta_k\pi\}$ ---
precisely because $\arg (w_k)$ is not uniquely determined.  However,
only one of these functions will be globally univalent on ${\cal D}$.
 In order to construct the univalent mapping $f$ from ${\cal D}$ to
${\cal G}$, i.e., in order to solve the accessory parameter problem
for the univalent mapping, we will need to measure the total changes
in argument over the circular sides of ${\cal G}$ rather than
measuring the absolute arguments at the vertices of ${\cal G}$.  We
will denote the change in argument over a circular side of ${\cal G}$
from vertex $w_{k-1}$ to $w_k$ by $\Delta \arg (w_k)$ and we will denote
the total change in argument for the mapping function $f$ over the
arc (on the boundary of ${\cal D}$) from $z_{k-1}$ to $z_k$ by $\Delta
\arg (w(k))$, where, again, $w(k) = f(z_k)$.

For each vertex $w_k , 1 \leq k \leq n-2,$ there are four (mutually
exclusive) possibilities that can arise: 
\begin{itemize}
\item[(a)] the vertex $w_k$  lies at infinity;
\item[(b)] the vertex $w_k$ is finite and lies on a radial side of
${\cal G}$ which comes in from infinity (as the boundary of ${\cal
G}$ is traversed positively), but not at the end of the radial side,
i.e., the interior angle at $w_k$ is $\pi$;  
\item[(c)] the vertex $w_k$ is finite and lies at the end of a radial
side of ${\cal G}$, which comes in from infinity;
\item[(d)] the vertex $w_k$ is otherwise finite.
\end{itemize}

If (a) holds, then we will not impose any condition on $w(k)$;
however, the geometry of ${\cal G}$ will force the next vertex to be
finite.

If (b) holds, then we will impose a modulus condition on $w(k)$ and
require that
\renewcommand{\theequation}{3a}
\begin{equation}
|w(k)| - |w_k| = 0.\label{eq:3a}\end{equation}

If (c) holds, then we will impose two conditions on $w(k)$.  We will
require that the modulus condition (\ref{eq:3a}) holds and we will
require that
\renewcommand{\theequation}{3b}
\begin{equation}
Arg (w(k)) - Arg (w_k) =0.\label{eq:3b}\end{equation}

In the last case, then either $|w_k| = |w_{k-1}|$ or else $Arg
(w_k) = Arg (w_{k-1})$.  If $|w_k| = |w_{k-1}|$, then we will
require
\renewcommand{\theequation}{3c}
\begin{equation}
\Delta \arg (w(k)) - \Delta \arg (w_k) = 0,\label{eq:3c}\end{equation}
\setcounter{equation}{3}otherwise we will require (\ref{eq:3a}) to
hold.

The above construction will generate $n-2$ equations if the vertex
$w_{n-2}$ does not lie at infinity.  If, however, $w_{n-2}$ does lie
at infinity, then one additional equation will need to be generated.
The vertex $w_{n-1}$ must be finite and satisfy condition
(c).  Thus, we can add one equation to the system by adding
the modulus equation (\ref{eq:3a}) with $k = n-1$.

The system of equations generated at this point may not be sufficient
to characterize the accessory parameter set for the correct mapping
function, especially in the case that ${\cal G}$ is unbounded and
non-starlike.  We may need to replace one of the equations in the
system with a secondary equation to control the total argument
change for the mapping function or to control the location of the tip
of a slit.  There are three cases which need to be considered.

Case 1.  The vertex $w_{n-1}$ or $w_n$ lies at the end of a radial
side of ${\cal G}$ going out to infinity and the interior angle at
that end is $\pi/2$.  If $w_{n-1}$ lies at the end of the radial
side, then the last equation in the system of type (\ref{eq:3b}) must
be replaced by an argument condition (\ref{eq:3c}) with $k = n-1$.
Alternately, if $w_n$ lies at the end of the radial side, then either
$w_{n-1}$ is the tip of a circular slit or the boundary of ${\cal G}$
has a right angle corner at $w_{n-1}$.  If $w_{n-1}$ is the tip of a
circular slit, then the last equation in the system of type
(\ref{eq:3a}) must be replaced by an argument condition (\ref{eq:3c})
with $k =n$; otherwise, the last equation in the system of type
(\ref{eq:3b}) must be replaced by an argument condition (\ref{eq:3c})
with $k=n$.

Case 2.  Case 1 does not apply and condition (c) holds for some
vertex $w_{k_0},\ 1 \leq k_0 \leq n-1$, such that $w_{k_0}$ lies on
the component of the boundary of ${\cal G}$ which contains $w_n$ and
the boundary of ${\cal G}$ has a right angle corner at $w_k$.  If
$w_{n-1}$ is the tip of a radial slit, then the argument equation
(\ref{eq:3b}) generated by condition (c) at $k = k_0$ must be
replaced by a modulus condition (\ref{eq:3a}) with $k = n-1$.
Alternatively, if $k_0 = n-1$, then the modulus equation (\ref{eq:3a})
generated by condition (c) at $k = k_0$ must be replaced by an
argument condition (\ref{eq:3c}) with $k=n$.  Alternately, if
$w_{n-1}$ lies at the tip of a circular slit, then the last equation
in the system of type (\ref{eq:3a}) must be replaced by an argument
condition (\ref{eq:3c}) with $k = n-1$.  Otherwise, if the segment on
the boundary of ${\cal G}$ between $w_{n-2}$ and $w_{n-1}$ is a
circular arc, then the argument equation (\ref{eq:3b}) generated by
condition (c) at $k = k_0$ must be replaced by an argument condition (\ref{eq:3c})
with $k = n-1$.

Case 3.  The domain ${\cal G}$ is bounded and the vertex $w_{n-1}$
lies at the tip of a slit.  The above construction does not
completely characterize the geometry of ${\cal G}$, i.e., the
generated mapping function $f$ may not satisfy $f(z_{n-1}) =
w_{n-1}$.  The last equation of the system, that is, the equation
representing the vertex $w_{n-2}$, must be replaced by a condition
determined by the geometry at the vertex $w_{n-1}$.  If $w_{n-1}$
lies at the tip of a radial slit, then we will require that
(\ref{eq:3a}) holds for $k = n-1$, otherwise we will require that
(\ref{eq:3c}) holds for $k = n-1$.

To determine the full set of $n-1$ accessory parameters we need to
add one additional equation to the above system of equations.  We
have not yet imposed the constraint condition (\ref{eq:iii}) in the
representation formula for gearlike mappings.  While we could solve
(\ref{eq:iii}) for one of the prevertices (in terms of the others)
and thus reduce the number of unknown accessory parameters from $n-1$
to $n-2$, our computational experience has been that we have been
able to treat a variety of problems (i.e., solve the accessory
parameter problem) using a full set of $n-1$ equations, with its
extra degree of freedom, which we could not successfully treat using
a reduced set of $n-2$ equations.  We will, therefore, impose
(\ref{eq:iii}) as the $n-1^{st}$ system equation.

Let us note that if the interior angle at $w_{n-1}$ is $\pi$,
then the generated mapping $f$ may not satisfy $f(z_{n-1}) = w_{n-1}$.
 No auxiliary condition can be added to the above construction
without over-determining the problem. On the other hand, the point
$w_{n-1}$ can be removed from the set of vertices of ${\cal G}$
without changing the geometry of ${\cal G}$.  We will require, in the
presentation of the vertices of ${\cal G}$, that $\beta_{n-1}$ be not
$0$.\vspace*{12pt}

\renewcommand{\theequation}{\arabic{equation}}\setcounter{equation}{3}
\noindent {\bf COMPUTATION}\vspace*{12pt}

In the actual practice of solving the nonlinear system of equations
(3) and computing the conformal mapping $f,$ we have generally
employed the techniques described by Trefethen.  (See \cite{10},[ pp.
86-90].)  The given form (3) requires solving for the unknown complex
prevertices $z_k$ on the unit circle $|z| = 1.$ (As with the
vertices, we will also denote $z_n$ as $z_0.$) It is numerically more
convenient to transform the points $z_k$ to their arguments
$\theta_k$ by 
$$z_k=e^{i\theta_k}, 0\leq\theta_k\leq 2\pi.$$ The
arguments $\theta_k$ are, however, severely constrained by a set of
linear inequalities because the points $z_k$ are ordered around the
unit circle.  The ordering constraints on (3) can be removed by
transforming the arguments $\theta_k$ to a set of unconstrained
variables $y\lower 5pt\hbox{$\scriptstyle k$}$ via the equations
\begin{equation}
y_k=\log\frac{\theta_k-\theta_{k-1}}{\theta_{k+1}-\theta_k}, \ 1\leq
k\leq n-1,\label{eq:4}\end{equation} where we take $\theta_0$ to be
$Arg(z_0) = 0$.  Because the arguments $\theta_k$ in (\ref{eq:4}) are
coupled, the transformation can be inverted to recover the arguments
$\theta_k$ from the images $y\lower 5pt\hbox{$\scriptstyle k$}$.

At each step in the iteration we will calculate from the
unconstrained variables $y_k , 1 \leq k \leq n-1,$ a set of arguments
$\theta_k$ and then a set of prevertices $z_k.$ (The prevertex $z_n$
is fixed at 1.)  Finally, we will calculate the values of the $n-1$
nonlinear equations (3) for the current set of prevertices.

In the computation of the values $f(z_k),$ we generally choose 
a path for the integration in (\ref{eq:2}) which is the straight line 
segment $[0,z_k ]$ in ${\cal D}.$  The integrand in (\ref{eq:2}) will have a
singularity  
at both $0$ and at $z_k.$  However, the singularity at $0$ is removable and 
can be controlled directly.  The singularity at $z_k$  is of the
form $(1 -\bar{z}_k z)^{-\beta_k}$  where  $\beta_k$  can take on only
one of five possible  
values.  Normally, such a singularity in an integral could be 
easily, and highly accurately, treated by Gauss-Jacobi quadrature.  
However, there may be, and typically there are, other prevertices 
$z_j$  clustered near $z_k,$ which could affect the accuracy of the 
quadrature result.  Trefethen described a type of compound Gauss-
Jacobi quadrature (see \cite{10}, [p. 87]) which divides the path of 
integration into subpaths, where the length of each subpath is 
dependent on how closely other prevertices $z_j$  are clustered to $z_k.$  
This type of compound Gauss-Jacobi quadrature has produced both 
highly accurate and efficient quadrature results.

We have used the library subroutine GAUSSQ by G. H. Golub and 
J. H. Welsch \cite{4} to calculate the nodes and weights for the 
Gauss-Jacobi quadrature.  We have usually set the number of nodes, 
NPTSQ, to be computed by GAUSSQ at 8.  Since the compound Gauss-
Jacobi quadrature always divides the path of integrations into two 
halves and one of the halves may be further subdivided, dependent 
on the distribution of the current set of prevertices $z_k,$ we 
achieve, in practice, at least 16 nodes for integration on each 
path.  

To solve the system of unconstrained nonlinear equations we 
have used the library subroutine NS01A by M. J. D. Powell \cite{8}, 
which employs a hybrid between the method of steepest descent 
(initially) and Newton's method (terminally).  

While the solution driver which calls NS01A can pass 
specialized initial sets of prevertices, should such information 
be available, the general initialization does not take into 
account any information about the geometry of the particular 
gearlike conformal mapping problem.  In general, the 
initialization allocates a set of initial prevertices uniformly 
distributed around the unit circle.  

Two of the control parameters required by NSO1A are: DSTEP, the step
size used by NS01A to calculate the Jacobian of (3) by
finite difference methods -- generally fixed at $10^{-8}$ and TOL,
the convergence criterion for returning from NS01A to the calling
program after having successfully solved (3).  The value of TOL is
closely related to the number of nodes NPTSQ selected for the
compound Gauss-Jacobi quadrature.  For a problem with vertices close
to unit norm TOL is typically set at $10^{{\rm -(NPTSQ+1)}}.$

Once the accessory parameters have been determined, values of the
mapping function $f$ can be computed for given initial points $z$ in
${\cal D}$ or on the boundary of ${\cal D}.$ In either case, the
values are computed by using compound Gauss-Jacobi quadrature.  While
the value of $f$ at a point $z$ can always be computed using (\ref{eq:2})
directly, if $z$ is near a point $z^* , z^* \neq 0,$ where the value
of $f$ is known, then the value $f(z)$ can be alternately computed as
$$f(z)=\frac{f(z^*)}{z^*}z\
\exp\left\{\int^z_{z^*}\frac{\left[\displaystyle{\prod^n_{k=1}}(1-\bar{z}_kw)^{-\beta_k}-1
\right]}{w}dw\right\}.$$ The point $z^*$ can be one of the prevertices
corresponding to a known finite vertex of ${\cal G}$ or the point
$z^*$ can be point where the value of $f$ has been previously
computed. 

Values for the inverse mapping function from the gearlike domain
${\cal G}$ to the unit disk ${\cal D}$ can also be generated.  Let
$w$ be in ${\cal G}$ or on the boundary of ${\cal G}.$ If, for $z =
f^{-1} (w),$ a nearby initial estimate $z^*$ can be given, then
Newton's method can be employed to solve for $z.$ (The value of the
derivative of $f$ can be calculated from equation (\ref{eq:1}).)  On
the other hand, when no initial estimate $z^*$ is known, we can
rewrite equation (\ref{eq:1}) as
\begin{equation}\frac{dw}{dz}=\frac{w}{z}\prod^n_{k=1}(1-\bar{z}_kz)^{-\beta_k}\label{eq:5}\end{equation}
where $w = f(z).$ The univalence of the mapping function implies that
equation
(\ref{eq:5}) can inverted to obtain
\begin{equation}\frac{dz}{dw}=\frac{z}{w}\prod^n_{k=1}(1-\bar{z}_kz).^{+\beta_k}\label{eq:6}\end{equation}
Equation (\ref{eq:6}) can be viewed as an ordinary differential
equation which lifts a straight line segment $[\tilde{w} ,w]$ in
${\cal G}$ to a solution curve in ${\cal D}$ from a point $\tilde{z}
= f^{-1} (\tilde{w})$ to the point $z.$ The point $\tilde{w}$ can be
any point in ${\cal G}$ provided that the straight line segment
$[\tilde{w} ,w]$ lies entirely in ${\cal G}.$ To solve the ordinary
differential equation we have used the library subroutine ODE by
Shampine and Gordon \cite{9}.  The code in ODE treats real differential
equations, which requires that we transform (\ref{eq:6}) to a coupled
pair of real differential equations.

We have, following Trefethen, combined the above two 
approaches, unless an initial estimate is known for which Newton's 
method can be applied directly.  We will first solve (\ref{eq:6}), starting 
typically at $\tilde{w}  = 0,$ to obtain a low order approximation $z^*$  to $z.$  
We will then follow up using Newton's method to move from $z^*$  to a 
final approximation of $z$ with high order accuracy.  

It is important that not only are we able to construct the 
approximate numerical mapping by solving the accessory parameter 
problem, but that we are able to estimate the accuracy of the 
approximation we have generated.  After the nonlinear solver NS01A 
returns a solution for the accessory parameter problem we check 
the accuracy of the computed solution by checking at each finite 
node $w_k$  the difference $|w_k  - w(k)|.$  At each infinite node $w_k$  we 
check the difference $\|w_{k-1}  - w_{k+1} | - |w(k-1) - w(k+1)\|.$   An 
error routine returns the maximum of the above tested 
differences as an error estimate for the accuracy of the generated 
mapping function.\vspace*{12pt}

\noindent {\bf EXAMPLES}\vspace*{12pt}
  
In Figure 3 we show conformal mappings generated at several of the
iterative steps which arise in constructing the conformal mapping in
Figure 2.  We start the process with no assumption about the geometry
of the prevertices for the image domain ${\cal G}$ in Figure 2, i.e.,
we begin with an initial uniform distribution about the unit circle
for the prevertices.  Also, we have set the number of nodes for the
compound Gauss-Jacobi quadrature at NPTSQ $= 8$ and have set the
level of desired accuracy at ${\rm TOL} = 10^{-8}.$ The plots in
Figure 3 show several contours (images of circles centered at the
origin) and the images of the six radial segments which join the
origin to the prevertices.  As Trefethen noted, the conformality of
the mappings at the origin can be used to interpolate on the plots in
Figure 3 the arguments of the prevertices via the arguments of the
images of the radial segments at they approach the origin.  The
example in Figure 3 is typical in that it shows some of the
prevertices to be clustering as the iterations move towards the
mapping solution at stage I20.\vspace*{12pt}

$\includegraphics[scale=0.22]{fig-03-a.eps}$ \hspace{0.25in} $\includegraphics[scale=0.22]{fig-03-b.eps}$
 
\noindent\hspace{1.25in}I5\hfill I7\hspace{1.25in}

\centerline{Figure 3}

\newpage

$\includegraphics[scale=0.22]{fig-03-c.eps}$ \hspace{0.25in} $\includegraphics[scale=0.22]{fig-03-d.eps}$
 
\noindent\hspace{1.25in}I8\hfill I9\hspace{1.25in}

\vspace*{12pt}

$\includegraphics[scale=0.22]{fig-03-e.eps}$ \hspace{0.25in} $\includegraphics[scale=0.22]{fig-03-f.eps}$
 
\noindent\hspace{1.25in}I11\hfill I20\hspace{1.25in}

\centerline{Figure 3 (cont.)}

\vspace*{12pt}

In Figures 4 and 5 we show examples which illustrate the 
dispersion of the contours for several sets of bounded and 
unbounded gearlike domains.  In each case the contours are the 
images of the circles $|z| = r,\ 0 < r < 1$.\vspace*{12pt}


In Figures 6 and 7 we show examples which illustrate 
the dispersion of the level sets for the inverse mappings for the 
gearlike domains depicted in Figures 4 and 5.  In each case the 
level sets are the images under $f^{-1}$  of $|w| = R,\ R >
0$.  The ticks shown on the boundary of ${\cal D}$ mark the location
of the prevertices $z_k$.


Finally, let us note one difficulty which we have encountered.
The only requirement for applying the above algorithm to a
gearlike domain mapping problem is that one of the finite vertices
must be distinguished as the last vertex, $w_n$, (and, or course,
then $w_{n-1}$ must not be removable).  There may be for a given
gearlike domain several vertices which could be so distinguished.
Our experience, however, in implementing the algorithm, i.e., in
solving the nonlinear system of parameter equations via NS01A, has
been that frequently only one of the several choices (for the
distinguished vertex) has led to a tractable system of equations.

\newpage

\vspace{20pt}

$\includegraphics[scale=0.22]{fig-04-a.eps}$ \hspace{0.25in} $\includegraphics[scale=0.22]{fig-04-b.eps}$
 
\noindent\hspace{1.25in}(a)\hfill (b)\hspace{1.25in}

\vspace*{12pt}

$\includegraphics[scale=0.22]{fig-04-c.eps}$ \hspace{0.25in} $\includegraphics[scale=0.22]{fig-04-d.eps}$
 
\noindent\hspace{1.25in}(c)\hfill (d)\hspace{1.25in}

\centerline{Figure 4}

\newpage

\vspace{20pt}

$\includegraphics[scale=0.22]{fig-05-a.eps}$ \hspace{0.25in} $\includegraphics[scale=0.22]{fig-05-b.eps}$
 
\noindent\hspace{1.25in}(a)\hfill (b)\hspace{1.25in}

\vspace*{12pt}

$\includegraphics[scale=0.22]{fig-05-c.eps}$ \hspace{0.25in} $\includegraphics[scale=0.22]{fig-05-d.eps}$
 
\noindent\hspace{1.25in}(c)\hfill (d)\hspace{1.25in}

\centerline{Figure 5}

\newpage

\vspace{20pt}

$\includegraphics[scale=0.22]{fig-06-a.eps}$ \hspace{0.25in} $\includegraphics[scale=0.22]{fig-06-b.eps}$
 
\noindent\hspace{1.25in}(a)\hfill (b)\hspace{1.25in}

\vspace*{12pt}

$\includegraphics[scale=0.22]{fig-06-c.eps}$ \hspace{0.25in} $\includegraphics[scale=0.22]{fig-06-d.eps}$
 
\noindent\hspace{1.25in}(c)\hfill (d)\hspace{1.25in}

\centerline{Figure 6}

\newpage

\vspace{20pt}

$\includegraphics[scale=0.22]{fig-07-a.eps}$ \hspace{0.25in} $\includegraphics[scale=0.22]{fig-07-b.eps}$
 
\noindent\hspace{1.25in}(a)\hfill (b)\hspace{1.25in}

$\includegraphics[scale=0.22]{fig-07-c.eps}$ \hspace{0.25in} $\includegraphics[scale=0.22]{fig-07-d.eps}$
 
\noindent\hspace{1.25in}(c)\hfill (d)\hspace{1.25in}

\centerline{Figure 7}

\newpage

\noindent {\bf APPLICATION}\vspace*{12pt}
   
Let $St$ be the class of analytic univalent functions $f$ on ${\cal
D},$ normalized by $f(0) = 0$ and $f'(0) = 1,$ which map ${\cal D}$
to an image domain which is starlike with respect to the origin.  For
$M \geq 1$ let ${\cal D} (M) = \{ z | \ |z| < M\}$ and let $St(M)$ be
the subclass of $St$ of functions which satisfy the property that
$f({\cal D} ) \subset {\cal D} (M).$ Let $f \in St (M)$ and have a
power series representation 
$$f(z)=z+\sum^\infty_{n=2}a_nz^n.$$
R. W. Barnard and J. L. Lewis \cite{1}, \cite{2} considered the problem of
finding \begin{equation}\begin{array}{cc} max & Re a_3\\ f\ \in\ St
(M)&\end{array}.\label{eq:7}\end{equation} They showed that if $M \geq
5$ and if $f$ is extremal for the coefficient problem (\ref{eq:7}),
then $f$ is the Pick function in $St (M),$ i.e., $f$ maps ${\cal D}$
onto the disk ${\cal D} (M)$ minus a radial slit on the negative real
axis.  (See Figure 8a.)  If $1 < M \leq e,$ they showed that if $f$
is extremal for (\ref{eq:7}), then $f$ is the SRT (square root
transform) Pick function in $St (M),$ i.e., $f$ maps ${\cal D}$ onto
the disk ${\cal D}(M)$ minus a pair of symmetric radial slits on the
imaginary axis.  (See Figure 8b.)\vspace*{12pt}

$$\includegraphics[scale=0.25]{fig-08.eps}$$

\noindent\hspace{1.2in}a\hfill b\hfill c\hspace{1.2in}

\centerline{Figure 8}
\vspace*{12pt}

While they were not able to solve the extremal problem for $e < M <
5,$ they were able to show that the extremal function must map ${\cal
D}$ onto ${\cal D} (M)$ minus a pair of real-axis symmetric radial
slits.  (See Figure 8c.)  They noted that if $M = 3,$ then the third
coefficient for the Pick function in $St (M)$ and the third
coefficient for the SRT Pick function in $St (M)$ are the same.  Lewis
posed the coefficient problem (\ref{eq:7}) for $e < M < 5$ as Problem
6.65 in \cite{3.5} and Barnard conjectured \cite{1} for $M \geq 3$ that
(\ref{eq:7}) was maximized by the Pick function in $St (M)$ and for $1 < M \leq 3$ that (\ref{eq:7}) was maximized by the SRT Pick
function in $St (M).$ 

To support that conjecture the problem was posed to show numerically
for $e < M < 5, M$ fixed, and for the real-axis symmetric double
radial slit mappings in $St (M)$ that $Rea_3$ was concave as a
function of the argument $\alpha$ of the omitted slit, $\pi /2 \leq
\alpha \leq \pi.$  Each such function $f$ is a bounded gearlike
mapping and can be represented by (from the geometry of the image
domain)
$$\frac{zf^\prime(z)}{f(z)} = \sqrt{\frac{p^2_2 (z)}{p_1 (z) p_3
(z)}} = Q(z)$$
where $p_j (z) = 1 - 2x_j z + z^2,\ x_j = Rez_j = \cos \theta_j,\ z_j
= e^{i\theta}j$ and $0 \leq \theta_1 \leq \theta_2 \leq \theta_3 \leq
 \pi$.  If we let
$$Q(z) = 1 + \sum^\infty_{k=1} q_k z^k$$
then
$$a_3 = \frac{q^2_1 + q_2}{2}$$
where 
\begin{eqnarray*}
q_1 & = & x_1 + x_3 - 2x_2\\[2ex]
q_2 & = & \frac{3(x_1 + x_3)^2 - 4(x_1 x_2 + x_2 x_3 + x_3
x_1)}{2}.\end{eqnarray*}


We specifically considered the case $M = 3.$ As
we constructed the values of $Rea_3$, for incremental values of
argument $\alpha$ of the omitted slit, we found (using a value for
$TOL = 10^{-8}$), however, (see Table
1) that the conjecture was not correct.  Indeed, Table 1 suggests
for $M=3$ that $Rea_3$ is unimodal as a function of $\alpha$ over the
interval $[\pi/2,\pi]$ and that the maximum value occurs for some
$\alpha$ in $[0.80\pi,0.90\pi]$.  

We considered various values for $M$ and found evidence to suggest
that the argument $\alpha$ of the omitted slit for the extremal
function for (\ref{eq:7}) appears to vary monotonically with $M$ from
$\pi /2$ to $\pi$ as $M$ varies from $e$ to $5.$

\vspace*{12pt}
\centerline{Table 1}
\vspace*{12pt}

\begin{center}
\begin{tabular}{ll}
\multicolumn{1}{c}{$\alpha/\pi$} & \multicolumn{1}{c}{$Re a_3$}\\\\
0.500 & 0.88888889\\
0.550 & 0.89011287\\
0.600 & 0.89358764\\
0.650 & 0.89873157\\
0.700 & 0.90461073\\
0.750 & 0.91000424\\
0.800 & 0.91353195\\
0.850 & 0.91388099\\
0.900 & 0.91009490\\
0.950 & 0.90174616\\
1.000 & 0.88888889\end{tabular}\end{center}\vspace*{12pt}

Barnard remarked, after seeing the numerical results in Table 1, that
it should be possible to show analytically that the conjecture was
false.  Indeed, a long, extensive calculus argument can be given
which shows that the conjecture fails for $2.83912 < M \leq 3$.
Specifically, the argument shows for $2.83912 < M$ that $Rea_3$ takes
on a local minimum value at the SRT Pick function over the set of
real-axis symmetric double radial slit mappings in $St(M)$.

We would like to thank Professor R. W. Barnard his helpful comments
and suggestions during the preparation of this paper. 

\newpage

\begin{thebibliography}{9999999999999}
\bibitem[Ba]{1} R. W. Barnard, ``A variational technique
for bounded starlike functions,'' {\em Can. J. Math.} 2 (75) 337-347.

\bibitem[Ba \& Le]{2} R. W. Barnard and J. L. Lewis,
``Coefficient bounds for some classes of starlike functions,'' {\em Pac. J.
Math.}  75   (1975), 325-331.

\bibitem[Ba \& Pe]{3} R. W. Barnard and K. Pearce, ``Rounding corners of 
gearlike domains and the omitted area problem,'' {\em J. Comp. 
Appl. Math.}  14  (1986), 217-226.

\bibitem[Bj \& Gr]{3.2} P. Bj\o rstad and E. Grosse, ``Conformal Mapping
of Circular Arc Polygons,'' SIAM J. Sci. Stat. Comp., to
appear.

\bibitem[Ca, Cl \& Ha]{3.5} D. M. Campbell, J. G. Clunie and W. K.
Hayman, ``Research problems in complex analysis,'' {\em Aspects of
contemporary complex analysis}, (Proc. NATO Adv. Study Inst., Univ.
Durham, 1979), 527-572, Academic Press, London, 1980.

\bibitem[Ga]{Ga} D. Gaier, Konstruktive Methoden der konformen
Abbildung, Springer Verlag, Berlin, 1964.

\bibitem[Go \& We]{4} 
G. H. Golub and J. H. Welsch, ``Calculation of Gaussian 
quadrature rules,'' {\em Math. Comp.}   23   (1969), 221-230.

\bibitem[Go]{5} A. W. Goodman, ``Conformal mapping onto certain
curvilinear polygons,'' Univ. Nac. Tucuman {\em A 13} (1960) 20-26.

\bibitem[He]{6} P. Henrici, {\em Applied and computational complex
analysis,} Vol. III, John Wiley \& Sons, New York, 1986.

\bibitem[Ja \& Ma]{6.5} K. P. Jackson and J. C. Mason, ``The
Approximation by Complex Functions of Stresses in Cracked Domains,''
{\em Algorithms for Approximation}, Oxford University Press, 1987.

\bibitem[Ne]{7} Z. Nehari, {\em Conformal Mapping}, McGraw-Hill Book Co.,
Inc., New York, 1952.

\bibitem[Po]{8} M. J. D. Powell, Technical Report AERE-R.5947, Harwell, 
England, 1968.

\bibitem[Sh \& Go]{9} L. Shampine and M. Gordon, {\em Computer
solution of ordinary differential equations: The initial value
problem}, Freeman, San Francisco, 1975.

\bibitem[Tr]{10} L. N. Trefethen, ``Numerical computation of the Schwarz-
Christoffel transformation,'' {\em SIAM J. Sci. Stat. Comp.}   1   
(1980), 82-102.

\bibitem[Tr, 83]{11} L. N. Trefethen, ``Radial and Gearlike
Domains,'' unpublished memo, 1983.

\end{thebibliography}
\end{document}


                 